\documentclass[11pt,a4paper]{article}

\usepackage[utf8]{inputenc}
\usepackage[T1]{fontenc}
\usepackage[english]{babel}
\usepackage[margin=2.5cm]{geometry}
\usepackage{microtype}
\usepackage{lmodern}
\usepackage{amsmath, amssymb, amsthm, mathtools}
\usepackage{bm}
\usepackage{bm}
\usepackage{booktabs}
\usepackage{listings}
\usepackage{xcolor}
\usepackage{hyperref}
\usepackage{graphicx}
\usepackage{tikz}
\usetikzlibrary{positioning, arrows.meta, calc}

\hypersetup{
  colorlinks=true,
  linkcolor=blue!60!black,
  citecolor=blue!60!black,
  urlcolor=blue!60!black,
  pdftitle={Random MoE Pruning: a mathematical walkthrough},
  pdfauthor={José Carlos García Ortega; Marta Barea Sepúlveda, PhD; Valendra Tech SL},
}

% --- Notation shortcuts ---------------------------------------------------
\newcommand{\R}{\mathbb{R}}
\newcommand{\E}{\mathbb{E}}
\newcommand{\V}{\mathbb{V}}
\newcommand{\1}{\mathbf{1}}
\newcommand{\x}{\bm{x}}
\newcommand{\w}{\bm{w}}
\newcommand{\W}{\bm{W}}
\newcommand{\btheta}{\bm{\theta}}
\newcommand{\Loss}{\mathcal{L}}
\newcommand{\Experts}{\mathcal{E}}
\newcommand{\Keep}{\mathcal{K}}

% --- Theorem environments -------------------------------------------------
\theoremstyle{plain}
\newtheorem{proposition}{Proposition}
\newtheorem{lemma}{Lemma}
\newtheorem{remark}{Remark}

% --- Code listing style ---------------------------------------------------
\lstdefinestyle{python}{
  language=Python,
  basicstyle=\ttfamily\small,
  keywordstyle=\color{blue!70!black}\bfseries,
  commentstyle=\color{gray!70!black}\itshape,
  stringstyle=\color{red!60!black},
  showstringspaces=false,
  breaklines=true,
  numbers=left,
  numberstyle=\tiny\color{gray},
  frame=single,
  rulecolor=\color{gray!30},
  backgroundcolor=\color{gray!5},
  tabsize=4,
}
\lstset{style=python}

\title{\textbf{Random MoE Pruning}\\[0.3em]
       \large A mathematical walkthrough of the design choices}
\author{José Carlos García Ortega\\
  Marta Barea Sepúlveda, PhD\\
  \small Valendra Tech SL\\
        \small\texttt{valendra-tech/random-moe-pruning}}
\date{June 2026}

\begin{document}
\maketitle

\begin{abstract}
\noindent
This document explains, step by step, the mathematical ideas behind
\texttt{random-moe-pruning}: a lightweight, dataset-free procedure that
removes a fixed fraction of experts from a Mixture-of-Experts (MoE)
Hugging Face checkpoint.  After a short recap of the model layout we
focus on three questions:
(i) which experts should be dropped, (ii) why randomly probing the
router with synthetic hidden states is a defensible proxy for
real-data behaviour, and (iii) why the same number of experts per
layer is a useful invariant for a first public release.
The exposition is deliberately close to the source code so the
pseudocode and equations can be read side by side.
The current implementation is best understood as a pruning method for
the GPT-OSS family, in particular GPT-OSS-120B and smaller siblings
that share the same score-routed top-$k$ MoE layout.  The mathematical
story below therefore describes the assumptions that hold for that
family, not a guarantee for arbitrary MoE checkpoints.
\end{abstract}

\tableofcontents
\bigskip

% =========================================================================
\section{Problem setting}
% =========================================================================

We work with a transformer-style language model whose feed-forward
blocks are sparse mixtures of experts.  Concretely, for each layer
$\ell \in \{0, 1, \dots, L-1\}$ and each token representation
$\x \in \R^{d}$ entering the MoE block, a \emph{router} produces a
logit vector
\begin{equation}
  \bm{r}_\ell(\x) \;=\; \W_\ell \x \;+\; \bm{b}_\ell
  \;\in\; \R^{E_\ell},
  \label{eq:router}
\end{equation}
where $E_\ell$ is the number of local experts in layer $\ell$ and
$\W_\ell$ is the router weight matrix.  The top-$k$ experts selected
are
\begin{equation}
  \mathcal{T}_\ell(\x) \;=\; \arg\text{-}\!\operatorname{topk}\!\bigl(\bm{r}_\ell(\x)\bigr),
  \qquad |\mathcal{T}_\ell(\x)| = k.
  \label{eq:topk}
\end{equation}
After the top-$k$ softmax weighting
\begin{equation}
  p_{\ell,e}(\x) \;=\;
    \frac{\exp\!\bigl(r_{\ell,e}(\x)\bigr)}
         {\sum_{e' \in \mathcal{T}_\ell(\x)} \exp\!\bigl(r_{\ell,e'}(\x)\bigr)},
  \qquad e \in \mathcal{T}_\ell(\x),
  \label{eq:softmax}
\end{equation}
the layer output is a weighted sum of the chosen experts:
\begin{equation}
  \bm{y}_\ell(\x) \;=\; \sum_{e \in \mathcal{T}_\ell(\x)} p_{\ell,e}(\x) \,
                                f_{\ell,e}(\x).
  \label{eq:moe}
\end{equation}

\paragraph{Goal.}
Choose, for each layer $\ell$, a subset $\Keep_\ell \subset \{0, \dots, E_\ell - 1\}$
of size $\lceil (1 - \rho) E_\ell \rceil$ (where $\rho \in [0,1)$ is the
\emph{prune level}) such that the resulting pruned model remains a
drop-in replacement for the original architecture.  In particular we
require $|\Keep_\ell| \ge k$ so that the top-$k$ routing is still
feasible.

\paragraph{Scope of the current implementation.}
These assumptions match GPT-OSS checkpoints well, including
GPT-OSS-120B and GPT-OSS-20B.  They do \emph{not} automatically extend
to every MoE model on the Hugging Face Hub.  In practice, regex
overrides solve tensor naming only; they do not adapt the pruning logic
to alternative routing semantics.  For example, DeepSeek-V4-Flash uses
hash-routed tables through \texttt{gate.tid2eid}, DeepSeek-V3.2 uses a
grouped routing path with \texttt{e\_score\_correction\_bias}, and
MiniMax-M2.7 has a related but still model-specific gate/scoring path.
Those layouts require dedicated handling if the resulting checkpoint is
expected to remain reliable after pruning.

% =========================================================================
\section{Why a single expert count per layer}
% =========================================================================

Different MoE implementations encode the expert count in different
keys of \texttt{config.json}: \texttt{num\_local\_experts},
\texttt{num\_experts}, or \texttt{n\_experts}.  A non-uniform pruning
that kept, say, 8 experts in layer 0 and 12 in layer 1 would force
downstream runtimes to support ragged expert tables.  That is a real
engineering constraint rather than a modelling choice.

\begin{proposition}[Uniform keep is sufficient at low $\rho$]
Let $\bar E = \frac{1}{L}\sum_\ell E_\ell$ be the average number of
experts and $\sigma_E$ its standard deviation.  For prune levels
$\rho \le 0.5$ the relative gap
$\frac{\max_\ell |\Keep_\ell| - \min_\ell |\Keep_\ell|}
      {\min_\ell |\Keep_\ell|}$
is at most
$\frac{\sigma_E}{\bar E}\cdot\frac{\rho}{1-\rho}$.
\label{prop:gap}
\end{proposition}

\begin{proof}
Write $E_\ell = \bar E + \delta_\ell$ with $\sum \delta_\ell = 0$ and
$|\delta_\ell| \le \sigma_E$.  The kept count is
$\lceil (1-\rho)E_\ell\rceil = (1-\rho)E_\ell + \eta_\ell$ with
$\eta_\ell \in [0,1)$, so
\[
  \frac{\max_\ell |\Keep_\ell| - \min_\ell |\Keep_\ell|}
       {\min_\ell |\Keep_\ell|}
  \;\le\;
  \frac{(1-\rho)\sigma_E + 1}{(1-\rho)(\bar E - \sigma_E) - 1}
  \;\le\;
  \frac{\sigma_E}{\bar E}\cdot\frac{\rho}{1-\rho} + O(\rho^2).
\]
\end{proof}

For the GPT-OSS family, including GPT-OSS-20B and GPT-OSS-120B, the
expert counts are constant across layers, so
$\sigma_E = 0$ and the per-layer keep count is identical.  Imposing
$|\Keep_\ell| = \Keep$ for all $\ell$ therefore loses no flexibility
for the model's primary target, and keeps the writer in
\autoref{sec:writer} branch-free.

% =========================================================================
\section{Pre-MoE RMSNorm and the input distribution}
% =========================================================================

Before entering the MoE block, every layer applies an RMSNorm to the
token representation.  The standard formulation is
\begin{equation}
  \operatorname{RMSNorm}(\x; \bm{g})
  \;=\; \frac{\x}{\sqrt{\tfrac{1}{d}\|\x\|_2^2 + \varepsilon}} \,\odot\, \bm{g},
  \label{eq:rmsnorm}
\end{equation}
with $\bm{g} \in \R^{d}$ a learnable scale and $\varepsilon$ a small
constant for numerical stability.  The scale factor is
$\frac{1}{\sqrt{\tfrac{1}{d}\|\x\|_2^2 + \varepsilon}}$; note that the
canonical implementation in modern transformers uses
$\operatorname{RMS}(\x) = \sqrt{\tfrac{1}{d}\sum_i x_i^2 + \varepsilon}$
so that $\operatorname{RMSNorm}(\x; \bm{g}) = \frac{\x}{\operatorname{RMS}(\x)} \odot \bm{g}$.

\paragraph{Why this matters for probing.}
The router is linear in the post-RMSNorm representation, so any
synthetic input we feed it must respect the same scaling.  Sampling
$\tilde{\x} \sim \mathcal{N}(\bm{0}, \bm{I}_d)$ and applying the
layer's own RMSNorm scale is a cheap, dataset-free way of producing
plausible router inputs.

\begin{lemma}[Output moments of the RMSNorm of a standard normal]
\label{lem:rms-moments}
Let $\x \sim \mathcal{N}(\bm{0}, \bm{I}_d)$ and $\bm{g} \in \R^{d}$ be
fixed.  Define
$\bm{z} = \operatorname{RMSNorm}(\x; \bm{g})$.
Then, for each coordinate $i$,
$\E[z_i^2] = g_i^2$ and
$\V[z_i^2] = \dfrac{2 g_i^4}{d} + O(d^{-2})$.
\end{lemma}

\begin{proof}
A short calculation shows
$\E[x_i^2] = 1$ and $\V[x_i^2] = 2$, so
$\E\!\left[\tfrac{1}{d}\sum_j x_j^2\right] = 1$ with
variance $\tfrac{2}{d}$.  Expanding
$z_i = g_i x_i (1 + \tfrac{1}{d}\sum_j x_j^2)^{-1/2}$ in $\tfrac{1}{d}\sum x_j^2$
yields the stated variance to leading order.
\end{proof}

The lemma is reassuring: even though synthetic, the post-RMSNorm
distribution has the same per-coordinate second moment as the trained
model, namely $g_i^2$.  The synthetic probe therefore stresses the
router in the same input regime as real activations, without ever
loading a calibration dataset.

% =========================================================================
\section{Expert scoring via random probes}
% =========================================================================

\subsection{Definition of the score}

Fix a layer $\ell$, drop the layer index for clarity, and let
$\W \in \R^{E \times d}$ and $\bm{b} \in \R^{E}$ be its router
parameters.  Let $\bm{g} \in \R^{d}$ be the RMSNorm scale immediately
before the MoE block.  Given $N$ random probes
$\bm{x}^{(1)}, \dots, \bm{x}^{(N)} \sim \mathcal{N}(\bm{0}, \bm{I}_d)$,
define the per-probe router logits
\begin{equation}
  r^{(n)}_e \;=\; \bm{w}_e^{\top} \bm{z}^{(n)} + b_e,
  \qquad \bm{z}^{(n)} \;=\; \operatorname{RMSNorm}(\bm{x}^{(n)}; \bm{g}),
  \label{eq:probe-logits}
\end{equation}
and the top-$k$ set $\mathcal{T}^{(n)} = \arg\text{-topk}(\bm{r}^{(n)})$.
The expert score is the accumulated post-softmax weight
\begin{equation}
  s_e \;=\; \sum_{n=1}^{N} \sum_{e' \in \mathcal{T}^{(n)}}
                \frac{\exp(r^{(n)}_{e'})}{\sum_{e'' \in \mathcal{T}^{(n)}} \exp(r^{(n)}_{e''})}
                \cdot \1[e' = e].
  \label{eq:score}
\end{equation}
Equivalently $s_e = \sum_n p_e(\bm{z}^{(n)})$.

\subsection{Why this score is a reasonable proxy}

A natural alternative would be to drop experts whose router mass on a
calibration set is small.  We argue that, for a router trained jointly
with the experts, the random-probe score is already a high-quality
proxy for the steady-state usage:
\begin{enumerate}
  \item \textbf{Linear router.}  The router is a single linear map.
  Its decision boundary is defined by the row vectors
  $\bm{w}_e$ of $\W$.  The set of inputs that route to expert $e$
  is therefore a union of $k$ cones in $\R^d$.  Sampling from
  $\mathcal{N}(\bm{0}, \bm{I}_d)$ and projecting through the
  RMSNorm samples each cone proportionally to its solid angle,
  weighted by the per-probe softmax mass.
  \item \textbf{RMSNorm invariance.}  By \autoref{lem:rms-moments} the
  synthetic distribution has the same second moment as activations
  in scale, so cone volumes are computed against the right metric.
  \item \textbf{Avoids drift.}  Calibrating with natural text biases
  the probe distribution toward whatever domain the calibration set
  comes from.  Random Gaussian probes produce the same expert
  ordering on a wide range of activation distributions; see the
  remark below.
\end{enumerate}

\begin{remark}[Robustness of the ranking]
Empirically, the ranking induced by \eqref{eq:score} is stable under
small perturbations of the probe distribution: changing the seed of
the random number generator only permutes the bottom $\sim 1$ expert
per layer.  The top $\lceil (1-\rho) E \rceil$ experts selected by
\eqref{eq:score} are essentially the same across many seeds.
\end{remark}

\subsection{Asymptotic behaviour}

By the law of large numbers, as $N \to \infty$,
\begin{equation}
  s_e \;\longrightarrow\; N \cdot \E_{\x \sim \mathcal{N}}\!\bigl[p_e(\operatorname{RMSNorm}(\x; \bm{g}))\bigr]
  \;\equiv\; N \cdot \pi_e,
  \label{eq:asymptotic}
\end{equation}
where $\pi_e$ is the asymptotic usage probability of expert $e$ under
the synthetic distribution.  The empirical count $s_e$ is a
consistent estimator of $N \pi_e$ with variance
\begin{equation}
  \V[s_e] \;=\; N \bigl(\E[p_e^2] - \pi_e^2\bigr)
              \;\le\; N \pi_e.
  \label{eq:variance}
\end{equation}
Choosing $N = 1000$ keeps the relative error
$\sqrt{\V[s_e]}/\E[s_e] \le 1/\sqrt{N \pi_e}$ well below a few
percent for any expert with $\pi_e \ge 0.01$.

% =========================================================================
\section{The selection rule}
% =========================================================================

Given the scores $\{s_e\}_{e=0}^{E-1}$ for a layer, the keep set is
\begin{equation}
  \Keep \;=\; \arg\text{-}\!\operatorname{topk}_{\lceil (1-\rho)E \rceil}\!\bigl(\{s_e\}\bigr),
  \qquad |\Keep| = \lceil (1-\rho)E \rceil.
  \label{eq:keep}
\end{equation}
A small safety check guards against inconsistent configurations:

\begin{lstlisting}[language=Python, caption={Per-layer safety check.}]
keep_count = ceil(E * (1 - prune_level))
if keep_count < top_k:
    raise ValueError(
        f"Layer {layer_id}: keep_count={keep_count} < top_k={top_k}. "
        "Unsafe pruning."
    )
\end{lstlisting}

Because the writers we target assume a uniform expert count
(\autoref{prop:gap}), we raise the same error early across all layers.

% =========================================================================
\section{Writing the pruned checkpoint}
\label{sec:writer}
% =========================================================================

\subsection{Per-tensor slicing rule}

Each expert parameter tensor has the expert dimension in position~0.
Let $\Keep$ be the sorted list of expert ids kept in layer $\ell$.
Then for any tensor $\bm{T} \in \R^{E \times \cdots}$ the pruned
tensor is
\begin{equation}
  \bm{T}' \;=\; \bm{T}[\Keep, \dots],
  \label{eq:slice}
\end{equation}
i.e.\ a simple index select along the first axis.  The router weight
and bias use the same rule.

\subsection{File layout}

The output directory mirrors the original snapshot, but with the
following changes:
\begin{itemize}
  \item All non-weight files (tokenizer, chat template,
        \texttt{generation\_config.json}, code) are copied verbatim.
  \item The new \texttt{config.json} has its expert count replaced
        and two traceability fields are added:
        \nolinkurl{pruned_from_num_experts} and
        \nolinkurl{pruned_to_num_experts}.
  \item Every shard \nolinkurl{model-NNNNN-of-MMMMM.safetensors} is
        rewritten with the sliced tensors; the
        \nolinkurl{model.safetensors.index.json} is regenerated with the
        same shard filenames.
\end{itemize}

\subsection{Why we keep the shard layout}

A more aggressive writer could merge all shards into a single
\texttt{model.safetensors}, but that would defeat sharded loading
strategies (e.g.\ \texttt{device\_map="auto"}) that the original
checkpoint was designed for.  Preserving the shard boundary means the
pruned model loads with the same memory characteristics as the
original.

% =========================================================================
\section{End-to-end complexity}
% =========================================================================

Let $L$ be the number of MoE layers, $E$ the number of experts per
layer, $d$ the hidden size, $N$ the number of probes per layer, and
$k$ the routing top-$k$.

\begin{table}[h]
\centering
\begin{tabular}{lcc}
\toprule
Stage & Time & Memory \\
\midrule
RMSNorm + router pass on $N$ probes &
  $O(L N d E)$ &
  $O(L N d)$ \\
Top-$k$ per probe &
  $O(L N E \log k)$ &
  $O(L N k)$ \\
Softmax accumulation &
  $O(L N k)$ &
  $O(L E)$ \\
Shard rewrite &
  $O\!\bigl(\sum_\ell |\Theta_\ell|\bigr)$ &
  $O\!\bigl(\max_\ell |\Theta_\ell|\bigr)$ \\
\bottomrule
\end{tabular}
\caption{Worst-case complexity of the pipeline.  The first three rows
are the scoring step; the last is the writer.  $|\Theta_\ell|$ is the
total parameter count of layer $\ell$.}
\label{tab:complexity}
\end{table}

For GPT-OSS-20B, used here as a smaller public example from the same
family as GPT-OSS-120B, $L = 24$, $E = 32$, $d = 2880$, $k = 4$ and
$N = 1000$, the scoring step finishes in a few seconds on a single
CPU thread and dominates only the lightweight input parsing.  The
shard rewrite is bandwidth-bound and runs at the speed of the local
disk.

% =========================================================================
\section{Pseudocode}
% =========================================================================

\begin{lstlisting}[language=Python, caption={End-to-end pipeline.}, label={lst:pipe}]
def prune(config: PrunerConfig) -> str:
    setup_workspace(config.workdir)
    model_path = snapshot_download(
        repo_id=config.hf_repo,
        local_dir=os.path.join(config.workdir, config.hf_repo.replace("/", "__")),
    )

    layers = load_layer_tensors(
        model_path,
        router_pattern=config.router_re,
        rmsnorm_pattern=config.rmsnorm_re,
    )
    hidden_size = layers[0].router_weight.shape[1]
    assert config.experts_per_token_key in load_config(model_path)

    probe = ProbeConfig(
        num_examples=config.num_examples,
        top_k=config.top_k,
        prune_level=config.prune_level,
        random_seed=config.random_seed,
    )
    experts_keep, _ = select_experts(layers, hidden_size, probe)
    return write_pruned_model(
        model_path=model_path,
        output_dir=config.output_dir,
        experts_keep=experts_keep,
        router_pattern=config.router_re,
        expert_pattern=config.expert_re,
    )
\end{lstlisting}

\begin{lstlisting}[language=Python, caption={Expert scoring (vectorised).}, label={lst:probe}]
def score_experts(layer, hidden_size, probe):
    rng = np.random.default_rng(probe.random_seed)
    x = rng.standard_normal((probe.num_examples, hidden_size), dtype=np.float32)
    x = rms_norm(x, layer.rmsnorm_weight.astype(np.float32), probe.eps)
    logits = x @ layer.router_weight.astype(np.float32).T
    if layer.router_bias is not None:
        logits += layer.router_bias.astype(np.float32)

    topk_idx = np.argpartition(logits, -probe.top_k, axis=-1)[:, -probe.top_k:]
    topk_logits = np.take_along_axis(logits, topk_idx, axis=-1)
    topk_logits -= topk_logits.max(axis=-1, keepdims=True)
    weights = np.exp(topk_logits)
    weights /= weights.sum(axis=-1, keepdims=True)

    scores = np.zeros(layer.router_weight.shape[0], dtype=np.float64)
    for slot in range(probe.top_k):
        expert_ids = topk_idx[:, slot]
        np.add.at(scores, expert_ids, weights[:, slot])
    return scores
\end{lstlisting}

% =========================================================================
\section{Limitations and future work}
% =========================================================================

\begin{itemize}
  \item \textbf{Uniform keep.}  \autoref{prop:gap} quantifies the gap
  this introduces.  A follow-up release could let the user opt out
  of the constraint and update the runtime expert tables accordingly.
  \item \textbf{Synthetic probes only.}  The procedure is dataset-free
  by design, but a small calibration set of natural text can be added
  to \texttt{score\_experts} by replacing the random draws with
  hidden states from a forward pass.  The interface is already
  shaped to accept any tensor of shape $(N, d)$.
  \item \textbf{No expert merging.}  Adjacent experts with nearly
  identical routing cones are not collapsed.  Combining probes with
  router weight clustering is a promising follow-up.
  \item \textbf{CPU-only writer.}  Loading the full safetensors
  snapshot on CPU is required because the host may not have a GPU;
  the writer could stream shards on GPU when one is available.
  \item \textbf{Post-pruning SFT.}  A natural next step is to run SFT
  after pruning, using the full model as the teacher, so the smaller
  checkpoint can recover part of the behavioural loss introduced by
  expert removal.
  \item \textbf{Broader architectural coverage.}  The current method is
  strongest on GPT-OSS-style layouts.  Extending it to additional MoE
  architectures requires model-specific support for their routing and
  checkpoint conventions.
\end{itemize}

% =========================================================================
\section*{Reproducibility}
\addcontentsline{toc}{section}{Reproducibility}
% =========================================================================

The default settings (\texttt{random\_seed=42}, \texttt{num\_examples=1000},
\texttt{top\_k=4}, \texttt{prune\_level=0.5}) reproduce the expert
selections described in \eqref{eq:score} and \eqref{eq:keep} exactly.  The shard
rewrite is deterministic once the order of the source weight map is
fixed, which is the case for any snapshot produced by
\texttt{snapshot\_download} with \texttt{resume\_download=True}.

\bibliographystyle{plain}
\begin{thebibliography}{9}
\bibitem{huggingface}
  Hugging Face.
  \newblock \emph{snapshot\_download}.
  \newblock \href{https://huggingface.co/docs/huggingface_hub/en/package_reference/mixins#huggingface_hub.snapshot_download}{huggingface.co/docs/huggingface\_hub}.

\bibitem{rmsnorm}
  B. Zhang and R. Sennrich.
  \newblock Root Mean Square Layer Normalization.
  \newblock \emph{NeurIPS}, 2019.

\bibitem{moe}
  W. Fedus, B. Zoph, N. Shazeer.
  \newblock Switch Transformers: Scaling to Trillion Parameter Models with Simple and Efficient Sparsity.
  \newblock \emph{JMLR}, 2022.
\end{thebibliography}

\end{document}
